\section{The composed ride-through protection level}
\label{sec:composition}

This is the object the paper is named for. It is a single protection level over
the whole degradation---multi-source cross-check, then handover, then coast---and
its one structural property is that it does not fall at the seam. The construction
and the property are \Proven{}; the specific ride-through magnitudes are
\Modelled{}.

\subsection{Definition}
\label{sec:compdef}

Let $\TPL_{\mathrm{handover}}$ be the scalar timing protection level of
Section~\ref{sec:scalar} evaluated with the source set present at the instant of
handover, and let $\HPL(\tau)$ be the $\Pzero$-seeded holdover envelope of
Section~\ref{sec:holdover}, with $\Pzero$ taken from the disciplining filter's
handover covariance. The composed ride-through protection level is
\begin{equation}
  \boxed{\;\PL(\tau)\;=\;\max\big(\TPL_{\mathrm{handover}},\,\HPL(\tau)\big),\quad \tau\ge0.\;}
  \label{eq:composed}
\end{equation}
For $\tau<0$ (before handover) the reported PL is the live $\TPL$ of the
cross-checked regime, whose value at the handover instant is
$\TPL_{\mathrm{handover}}=\PL(0^{-})$.

\subsection{Lower-semicontinuity across handover}
\label{sec:lsc}

The naive expectation is that a correctly seeded holdover PL should be
\emph{continuous} at the seam---that $\HPL(0)=\TPL_{\mathrm{handover}}$. That
expectation is false, and acting on it is unsafe. The handover $\TPL$ carries the
fault-mode overbounds of the integrity monitor: it protects against the
possibility that one of the several sources was faulty up to its minimum detectable
bias. A pure coast has no such fault hypotheses---there are no sources left to be
faulty---so a freshly re-derived $\HPL(0)$ built only from the clock's phase
variance is generically \emph{smaller} than $\TPL_{\mathrm{handover}}$. Demanding
continuity would then pull the PL \emph{down} at $\tau=0$, understating the
integrity risk exactly when the last independent check is lost.

\begin{theorem}[Lower-semicontinuity and non-decrease at handover]
\label{thm:lsc}
For every handover $\TPL_{\mathrm{handover}}\ge0$ and every holdover envelope
$\HPL(\cdot)\ge0$, the composed PL \eqref{eq:composed} satisfies
\begin{equation}
  \PL(0^{+}) \;=\; \max\big(\TPL_{\mathrm{handover}},\,\HPL(0)\big)
  \;\ge\; \TPL_{\mathrm{handover}} \;=\; \PL(0^{-}),
  \label{eq:lsc}
\end{equation}
so the composed PL is lower-semicontinuous and non-decreasing across handover: it
never jumps downward at the seam. Moreover, on any interval where $\HPL$ is
non-decreasing (Section~\ref{sec:holdover}), $\PL(\tau)$ is non-decreasing on
$\tau\ge0$.
\end{theorem}

\begin{proof}
Immediate from $\max(a,b)\ge a$: the first equality is the definition and the
inequality is the max property, giving \eqref{eq:lsc}. Non-decrease on
$\tau\ge0$ follows because the pointwise maximum of the constant
$\TPL_{\mathrm{handover}}$ and a non-decreasing $\HPL(\tau)$ is non-decreasing.
\end{proof}

The mathematical content of Theorem~\ref{thm:lsc} is deliberately elementary---it is
the one-line observation $\max(a,b)\ge a$. The contribution is not that algebra but
the \emph{modelling insight} it formalizes: that demanding continuity at the
cross-check--to--holdover seam is unsafe (because $\TPL_{\mathrm{handover}}$ carries
fault-mode overbounds a coast lacks), and that the $\max(\cdot)$ is the remedy that
guarantees no downward jump. ``Lower-semicontinuity'' is the \emph{name} of the
resulting property, not the achievement. We are precise about what this theorem is
and is not. It is \emph{not} a continuity claim and \emph{not}
$\HPL(0)=\TPL_{\mathrm{handover}}$; it is the weaker, correct statement that the PL
does not drop. The engine's test
\code{pl\_is\_lower\_semicontinuous\_across\_handover} asserts
$\PL(0^{+})\ge\TPL_{\mathrm{handover}}$ for a range of handover TPLs, and
\code{composed\_pl\_is\_non\_decreasing\_over\_coast} asserts monotonicity along
the coast. Theorem~\ref{thm:lsc} is the fix for a genuine defect: an earlier,
naive ``continuous and monotone at handover'' formulation is \emph{wrong} because
it forces the downward jump; the $\max(\cdot)$ construction is what makes the seam
safe.

\subsection{The $\Pzero$ seed is a safety requirement}
\label{sec:p0safety}

Theorem~\ref{thm:lsc} holds structurally through the $\max(\cdot)$, but the seed
$\Pzero$ still matters, for a distinct reason. Consider the raw holdover branch
\emph{without} the composition---the PL a receiver would report if it simply
switched to $\HPL(\tau)$ at handover. With $\Pzero=0$ (coasting as though the clock
were perfectly known) and no aging, $\HPL(0)=0$: the raw holdover PL starts at
zero and lies below any positive $\TPL_{\mathrm{handover}}$, an actively unsafe
drop. The engine's test
\code{unseeded\_raw\_hpl\_would\_drop\_but\_composed\_pl\_does\_not} exhibits both
halves: the raw $\HPL(0)$ with $\Pzero=0$ falls below a
\SI{30}{\nano\second} handover TPL, while the composed PL \eqref{eq:composed} holds
at the TPL. Seeding $\Pzero$ from the handover covariance is therefore a safety
requirement in its own right---it makes the holdover branch itself start above
zero---and the $\max(\cdot)$ is the belt-and-braces guarantee that even a
mis-seeded branch cannot pull the reported PL below the handover value.

\subsection{A representative ride-through}
\label{sec:ridethrough}

Table~\ref{tab:ridethrough} shows the composed PL for a representative
three-source handover (a GNSS-PPP source $\{\sigma=\SI{2}{\nano\second},
b=\SI{1}{\nano\second}\}$, a PTP/White-Rabbit source $\{\SI{5}{\nano\second},
\SI{2}{\nano\second}\}$, and a LEO-STL source $\{\SI{8}{\nano\second},
\SI{3}{\nano\second}\}$, each $p_{\mathrm{fault}}=\num{e-4}$) at integrity-risk
budget $\IR=\num{e-6}$ and per-mode false-alert budget $p_{\mathrm{fa}}=\num{e-3}$,
coasting on a disciplined oscillator. The handover TPL is
$\TPL_{\mathrm{handover}}=\SI{26.961}{\nano\second}$, driven (Theorem
\ref{thm:biasdriven}) by the exclusion of source~0. The composed PL stays flat at
this value while $\HPL(\tau)<\TPL_{\mathrm{handover}}$---through the first hour of
coast---then rises as the holdover envelope overtakes it, reaching
\SI{264.768}{\nano\second} at one day and \SI{32185.028}{\nano\second} at one week,
where the $\tau^{5}$ drift term dominates. Every number is regenerated by the
committed example (Section~\ref{sec:availability}); these are \Modelled{}
magnitudes for a representative clock, carrying no certification implication.

\begin{table}[t]
  \centering
  \caption{Composed ride-through PL $\PL(\tau)=\max(\TPL_{\mathrm{handover}},
  \HPL(\tau))$ for the representative three-source handover
  ($\IR=\num{e-6}$, $p_{\mathrm{fa}}=\num{e-3}$). The PL is flat at the handover
  TPL through the first hour, then holdover-dominated. \Modelled{}.}
  \label{tab:ridethrough}
  \begin{tabular}{r r l}
    \toprule
    coast $\tau$ & $\PL(\tau)$ (ns) & regime \\
    \midrule
    \SI{0}{\second}        & 26.961      & handover TPL \\
    \SI{60}{\second}       & 26.961      & TPL-limited \\
    \SI{600}{\second}      & 26.961      & TPL-limited \\
    \SI{3600}{\second}     & 26.961      & TPL-limited \\
    \SI{86400}{\second} (1\,d)  & 264.768     & holdover-dominated \\
    \SI{604800}{\second} (1\,wk) & 32185.028   & holdover-dominated ($\tau^{5}$) \\
    \bottomrule
  \end{tabular}
\end{table}

\subsection{Relation to the single-source conditional TPL}
\label{sec:baweja}

Baweja's conditional timing protection level \citep{baweja2026tpl} composes a
detector floor with a coast-over-latency term for a \emph{single}
GNSS-disciplined source across \emph{one} transition. It is the $N=1$, one-seam
special case of \eqref{eq:composed}: its detector floor is the integrity-risk-budgeted
counterpart (analogue) of the $N=1$ scalar TPL \eqref{eq:n1}---Baweja's floor is a
fixed-$k$ cross-satellite monitor floor, not the exclusion-free $N=1$ MHSS quantile,
so allocating the integrity-risk budget is itself part of the generalization---and
its coast term is the $\HPL(\tau)$ leg. The present construction generalizes it in
two ways. First, to $N\ge1$ multi-source operation through the rank-one MHSS budget,
taking the per-source overbounds $(\sigma_i,b_i)$ as \emph{given inputs}:
``heterogeneous'' here means multi-source with given overbounds, \emph{not} the full
$\utck{}$-traceability treatment---the UTC$(k)$-bias library and the correlated
cross-covariance are deferred to the companion paper of
Section~\ref{sec:sourcemodel}, and independence across sources is assumed here
(Section~\ref{sec:sourcemodel}). Second, to a lower-semicontinuous composition whose
safety at the seam is a proved property rather than a single-source construction. We
claim the generalization, and we cite the special case; we do not restate it as new.
